\documentclass[11pt]{article}
\usepackage[margin=1in]{geometry}
\usepackage[utf8]{inputenc}
\usepackage[T1]{fontenc}
\usepackage{lmodern}
\usepackage{natbib}
\usepackage{hyperref}
\usepackage{url}
\usepackage{booktabs}
\usepackage{amsfonts}
\usepackage{nicefrac}
\usepackage{microtype}
\usepackage{graphicx}
\graphicspath{{figures/}}
\usepackage{xcolor}
\usepackage{amsmath}
\usepackage{multirow}
\usepackage{subcaption}

\title{Do Reasoning Models Ask Better Questions?\\A Formal Information-Theoretic Analysis of Multi-Turn LLM Games}

\author{%
  Daniel Machado\thanks{Equal contribution.} \\
  AKCIT / UFG \\
  Goiânia, Brazil \\
  \texttt{daniel@akcit.ai} \\
  \And
  Bryan L. M. Oliveira\footnotemark[1] \\
  AKCIT / UFG \\
  Goiânia, Brazil \\
  \texttt{bryan@akcit.ai} \\
  \And
  % Add other authors here
}

\begin{document}

\maketitle

\begin{abstract}
Large Language Models excel at many tasks but struggle with a critical capability for autonomous agents: asking informative questions to resolve ambiguity under uncertainty. We introduce \textbf{InfoGainME}, an information-theoretic framework for evaluating LLMs' question-asking ability through multi-turn yes/no dialogue games over structured candidate pools. Unlike prior work using external scaffolding or preference optimization, we study whether \emph{intrinsic} Chain-of-Thought (CoT) reasoning enables models to implicitly perform Bayesian active learning. Across 25,513 games spanning three domains (geography, objects, diseases), seven model sizes (0.6B--235B), and two observability conditions, we find that: (1) CoT provides substantial benefits for models $\geq$4B parameters, with the largest effect (+0.30 win rate) in partially observable settings; (2) below 4B parameters, CoT is actively harmful; (3) larger models shift from attribute-based to feature-based questioning, leveraging world knowledge beyond explicit metadata; (4) 46\% of all turns produce zero information gain, with 43\% of these occurring when only one candidate remains---revealing a systematic failure to recognize convergence; and (5) CoT eliminates question redundancy, suggesting reasoning traces serve as implicit dialogue memory. Our analysis provides the first rigorous decomposition of how reasoning models navigate hypothesis spaces under uncertainty.
\end{abstract}

\section{Introduction}
\label{sec:intro}

Information-seeking is fundamental to intelligent behavior. When an autonomous agent faces ambiguity---whether diagnosing a medical condition, troubleshooting a system, or understanding user intent---it must identify what information is missing and ask questions that efficiently reduce uncertainty. This capability, formalized through expected information gain (EIG), is central to Bayesian experimental design \citep{lindley1956measure,chaloner1995bayesian} and has recently attracted attention in the context of Large Language Models (LLMs).

Despite impressive performance on reasoning benchmarks, LLMs struggle with proactive information-seeking. Prior work has documented that vanilla LLMs ask uninformative questions \citep{bertolazzi2023lets}, fail to update beliefs consistently across dialogue turns \citep{choudhury2025bedllm}, and require external scaffolding or fine-tuning to improve \citep{mazzaccara2024learning,hu2024uncertainty}. However, the emergence of \emph{reasoning models}---LLMs trained or prompted to generate explicit Chain-of-Thought (CoT) traces before answering---raises a fundamental question:

\begin{center}
\emph{Do reasoning traces enable LLMs to implicitly perform optimal information-seeking?}
\end{center}

This question is orthogonal to prior work. Approaches like BED-LLM \citep{choudhury2025bedllm} use external hypothesis enumeration and EIG computation at deployment time, while DPO-based methods \citep{mazzaccara2024learning} fine-tune models on high-EIG question pairs. Neither asks whether \emph{intrinsic} reasoning---the deliberation visible in CoT traces---allows models to approximate Bayesian active learning without external intervention.

We introduce \textbf{InfoGainME} (Information Gain in Multi-turn Evaluation), a framework for studying this question through controlled dialogue games. Our contributions are:

\begin{enumerate}
    \item \textbf{A rigorous evaluation framework.} InfoGainME uses structured candidate pools with rich metadata (160 cities, 158 objects, 160 diseases), enabling precise measurement of information gain per turn, comparison to the theoretical optimal baseline (1 bit/turn), and cross-domain generalization analysis.
    
    \item \textbf{The first systematic study of reasoning models in information-seeking.} We compare CoT vs.\ non-CoT variants across seven model sizes (0.6B--235B), isolating the effect of reasoning from raw scale.
    
    \item \textbf{Reasoning trace analysis.} We operationalize the analysis of CoT traces through candidate question diversity, question-type taxonomy (attribute-based, feature-based, direct guess), and redundancy detection---providing the first quantitative decomposition of \emph{how} reasoning models explore hypothesis spaces.
    
    \item \textbf{A comprehensive failure analysis.} We categorize 167,767 zero-IG turns into root causes, discovering that 43\% occur when only one candidate remains---a systematic failure to recognize convergence and assert the answer.
\end{enumerate}

\section{Related Work}
\label{sec:related}

\paragraph{Information-Seeking Dialogue.}
Traditional approaches to clarification and disambiguation include benchmarks like ClariQ \citep{aliannejadi2021building}, QReCC \citep{anantha2021open}, and AmbigQA \citep{min2020ambigqa}. These focus on generating clarifying questions but typically lack information-theoretic evaluation. \citet{deng2023prompting} showed that open-source LLMs struggle with proactive clarification, predominantly reacting to queries rather than taking initiative.

\paragraph{20-Questions with LLMs.}
The 20-questions game provides a natural testbed for information-theoretic question quality. \citet{bertolazzi2023lets} evaluated ChatGPT using hierarchical hypothesis spaces, finding that LLM questions approach optimality only when \emph{explicitly prompted} to enumerate updated hypotheses step-by-step. \citet{mazzaccara2024learning} improved Llama-2-7b using Direct Preference Optimization (DPO) on high-EIG vs.\ low-EIG question pairs, demonstrating that EIG serves as an effective training signal. \citet{zhang2024eda} introduced the Entity-Deduction Arena (EDA) benchmark, showing that GPT-4 outperforms humans and that RL training substantially improves weaker models. LMRL-Gym \citep{abdulhai2023lmrl} established that standard LLMs fail at multi-turn information gathering, motivating RL-based approaches. Paprika \citep{tajwar2025paprika} trains generally curious agents that transfer exploration skills across environments.

\paragraph{Bayesian Experimental Design with LLMs.}
BED-LLM \citep{choudhury2025bedllm} proposes sequential Bayesian experimental design for LLM information gathering, explicitly enumerating hypothesis pools and selecting questions to maximize EIG. IGPO \citep{wang2025igpo} uses turn-level information gain as an intrinsic reward for RL training. Both demonstrate that EIG is an effective optimization objective; our work uses EIG as an \emph{evaluation} metric to compare model families.

\paragraph{Uncertainty and Planning in LLMs.}
Uncertainty of Thoughts (UoT) \citep{hu2024uncertainty} augments LLMs with uncertainty-aware planning, achieving 38.1\% improvement on information-seeking tasks via explicit EIG maximization at inference time. Tree of Thoughts \citep{yao2023tot} enables deliberate search over reasoning paths with backtracking. While these approaches \emph{augment} models at inference time, we ask whether \emph{reasoning models}---trained via RLVR to generate CoT traces \citep{deepseek2025r1,openai2024o1,qwen2025qwen3}---have internalized similar capabilities without external scaffolding.

\paragraph{Chain-of-Thought and its Limitations.}
\citet{wei2022chain} established that CoT prompting improves complex reasoning in large models. However, \citet{sprague2024cot} found that CoT benefits primarily math and symbolic tasks, with smaller gains elsewhere. \citet{stechly2024thoughtlessness} documented cascading errors in CoT planning. Our work examines a domain not covered in prior CoT analyses: interactive, belief-tracking, multi-turn information gathering---where we hypothesize reasoning models' explicit deliberation should provide benefits analogous to implicit belief state maintenance.

\paragraph{Multi-Turn Failures.}
\citet{laban2025lost} found that all top LLMs exhibit 39\% performance drops in multi-turn vs.\ single-turn settings, with a key failure mode being premature commitment to early assumptions. Our failure analysis (Section~\ref{sec:results}) provides a formal decomposition of multi-turn failures through the information gain lens.

\section{The InfoGainME Framework}
\label{sec:framework}

\subsection{Game Structure}

InfoGainME implements a multi-turn dialogue game with three roles:

\begin{itemize}
    \item \textbf{Seeker}: The evaluated LLM, which asks yes/no questions to identify a hidden target from a candidate pool.
    \item \textbf{Oracle}: Knows the target and provides truthful yes/no answers.
    \item \textbf{Pruner}: Eliminates candidates inconsistent with each Oracle response, maintaining the active candidate set.
\end{itemize}

The Seeker operates in one of two observability modes:
\begin{itemize}
    \item \textbf{Fully Observable (FO)}: Sees the complete list of active candidates with metadata at each turn.
    \item \textbf{Partially Observable (PO)}: Sees only the dialogue history; must implicitly track remaining candidates.
\end{itemize}

\subsection{Domains and Candidate Pools}

We evaluate across three domains with distinct metadata structures:

\begin{itemize}
    \item \textbf{Geography}: 160 cities from the world's most populous, annotated with region, subregion, country, and state.
    \item \textbf{Objects}: 158 everyday items (animals, vehicles, household objects, etc.), each with a category label.
    \item \textbf{Diseases}: 160 medical conditions, each annotated with associated symptoms.
\end{itemize}

The structured metadata enables both targeted attribute-based questioning and measurement of strategy quality.

\subsection{Information Gain Metric}

We measure question quality through Shannon information gain:

\begin{equation}
    \text{IG} = H_{\text{before}} - H_{\text{after}} = \log_2(N_{\text{before}}) - \log_2(N_{\text{after}})
\end{equation}

where $N$ is the number of active candidates. Under a uniform prior, the theoretical optimum is 1 bit/turn (perfect binary search halves the candidate pool each turn). For a pool of size $N$, the optimal number of turns is $\lceil \log_2 N \rceil \approx 8$ for our pools.

\section{Experimental Setup}
\label{sec:setup}

\subsection{Models}

We evaluate models from the Qwen3 family spanning five orders of magnitude:
\begin{itemize}
    \item \textbf{0.6B}: Qwen3-0.6B
    \item \textbf{4B}: Qwen3-4B-Instruct, Qwen3-4B-Thinking
    \item \textbf{8B}: Qwen3-8B (with and without CoT prompting)
    \item \textbf{30B}: Qwen3-30B-A3B-Instruct, Qwen3-30B-A3B-Thinking
    \item \textbf{235B}: Qwen3-235B-A22B-Instruct (FP8)
\end{itemize}

The Oracle and Pruner use Qwen3-8B by default, with Gemma-3-12B-IT for validation.

\subsection{Experimental Conditions}

Each model is tested under:
\begin{itemize}
    \item \textbf{Observability}: FO vs.\ PO
    \item \textbf{Reasoning}: CoT-enabled (Thinking variants or CoT prompting) vs.\ standard instruction-tuned
    \item \textbf{Domains}: Geography, Objects, Diseases
\end{itemize}

We run 3 trials per target per condition, yielding 25,513 total conversations and 363,846 turns.

\section{Results}
\label{sec:results}

\subsection{RQ1: Do Reasoning Models Optimize for Information Gain?}

Figure~\ref{fig:cot_benefit} shows the effect of CoT on win rate across model sizes and observability modes. Table~\ref{tab:rq1} provides detailed metrics.

\begin{figure}[h]
\centering
\includegraphics[width=0.85\textwidth]{fig1_cot_benefit.pdf}
\caption{Effect of Chain-of-Thought on win rate. CoT is harmful for 0.6B (red region) but beneficial for $\geq$4B (green region), with the largest effect in partially observable settings.}
\label{fig:cot_benefit}
\end{figure}

Table~\ref{tab:rq1} compares CoT vs.\ non-CoT variants across model sizes.

\begin{table}[h]
\centering
\caption{Win rate and IG/turn by model size, observability, and reasoning mode. Bold indicates CoT improvement $>$0.05.}
\label{tab:rq1}
\small
\begin{tabular}{llcccccc}
\toprule
Size & Mode & CoT WR & no-CoT WR & $\Delta$ WR & CoT IG/t & no-CoT IG/t & $\Delta$ IG/t \\
\midrule
0.6B & FO & 0.833 & 0.968 & -0.135 & 2.969 & 3.628 & -0.658 \\
0.6B & PO & 0.713 & 0.900 & -0.187 & 3.380 & 4.261 & -0.880 \\
\midrule
4B & FO & 0.978 & 0.892 & \textbf{+0.086} & 0.796 & 0.691 & \textbf{+0.105} \\
4B & PO & 0.574 & 0.277 & \textbf{+0.297} & 0.428 & 0.302 & \textbf{+0.126} \\
\midrule
8B & FO & 0.999 & 0.973 & +0.026 & 0.797 & 0.776 & +0.021 \\
8B & PO & 0.626 & 0.598 & +0.028 & 0.443 & 0.430 & +0.013 \\
\midrule
30B & FO & 0.994 & 0.999 & -0.006 & 0.873 & 0.805 & \textbf{+0.068} \\
30B & PO & 0.635 & 0.511 & \textbf{+0.124} & 0.463 & 0.378 & \textbf{+0.086} \\
\bottomrule
\end{tabular}
\end{table}

\textbf{Key findings:}
\begin{itemize}
    \item CoT provides substantial benefits for models $\geq$4B, with the largest effect in PO mode. The 4B model shows +0.30 win rate improvement in PO---effectively doubling performance from 28\% to 57\%.
    \item \textbf{CoT is harmful for 0.6B models} (-0.14 to -0.19 win rate). The high IG/turn values for 0.6B are artifacts of aggressive early guessing, not better information-seeking.
    \item In FO mode, models $\geq$8B already achieve $>$97\% win rate, creating a ceiling effect. The CoT benefit is small but positive for IG efficiency.
    \item \textbf{PO is where CoT matters}: Consistent improvement across 4B/8B/30B confirms that reasoning traces help maintain implicit belief states.
\end{itemize}

\subsection{RQ2: Does Scale Matter?}

In the hardest condition (PO + CoT):
\begin{itemize}
    \item 4B: 57\% $\rightarrow$ 8B: 63\% $\rightarrow$ 30B: 64\%
\end{itemize}

Size helps with diminishing returns. The dramatic improvement from 0.6B to 4B (not shown in CoT metrics due to CoT being harmful at 0.6B) suggests \textbf{$\sim$4B parameters is the minimum threshold for effective reasoning-augmented information seeking}.

\subsection{RQ3: How Do Reasoning Models Explore?}

\paragraph{Question Type Analysis.}
We classified 363,846 questions into attribute-based (targeting explicit metadata), feature-based (using world knowledge beyond metadata), and direct guesses. Figure~\ref{fig:question_types} shows the distribution across model configurations.

\begin{figure}[h]
\centering
\includegraphics[width=0.9\textwidth]{fig2_question_types.pdf}
\caption{Question type distribution by model configuration. Larger models (30B+) shift from attribute-based to feature-based questioning, leveraging world knowledge beyond explicit metadata. The 235B model in PO asks 65\% feature-based questions.}
\label{fig:question_types}
\end{figure}

\begin{table}[h]
\centering
\caption{Question type distribution by model configuration.}
\label{tab:question_types}
\small
\begin{tabular}{lcccc}
\toprule
Config & Attribute & Feature & Direct & Other \\
\midrule
4B FO no-CoT & 71.0\% & 23.0\% & 6.0\% & 0.0\% \\
8B FO CoT & 68.6\% & 23.4\% & 8.1\% & 0.0\% \\
30B FO Thinking & 37.7\% & 52.1\% & 10.2\% & 0.0\% \\
235B PO no-CoT & 34.7\% & 65.3\% & 0.0\% & 0.0\% \\
0.6B FO CoT & 61.6\% & 27.4\% & 2.1\% & 8.9\% \\
\bottomrule
\end{tabular}
\end{table}

\textbf{Key finding}: Larger models (30B+) shift from attribute-based to feature-based questions. The 235B model in PO asks 65\% feature-based questions---it uses world knowledge to generate discriminative queries \emph{beyond} the explicit metadata. This represents emergent Bayesian reasoning through world knowledge.

\paragraph{Redundancy Analysis.}
CoT dramatically reduces question redundancy (Figure~\ref{fig:redundancy}):
\begin{itemize}
    \item 4B-Instruct: 5.2\% $\rightarrow$ 0\% with CoT
    \item 30B-Instruct: 5.7\% $\rightarrow$ 0\% with CoT
\end{itemize}

\begin{figure}[h]
\centering
\includegraphics[width=0.7\textwidth]{fig7_redundancy.pdf}
\caption{CoT eliminates question redundancy across model sizes. Check marks indicate models that achieve 0\% redundancy with CoT.}
\label{fig:redundancy}
\end{figure}

This suggests reasoning traces serve as implicit dialogue memory, preventing models from re-asking equivalent questions.

\subsection{RQ4: Performance vs.\ Optimal}

Figure~\ref{fig:efficiency} shows efficiency relative to the theoretical optimal (1 bit/turn) across domains.

\begin{figure}[h]
\centering
\includegraphics[width=0.8\textwidth]{fig5_efficiency.pdf}
\caption{Information gathering efficiency relative to theoretical optimal (perfect binary search). Even in the best case (geography FO), models achieve only 80\% efficiency.}
\label{fig:efficiency}
\end{figure}

\begin{table}[h]
\centering
\caption{Efficiency relative to theoretical optimal (1 bit/turn).}
\label{tab:optimal}
\small
\begin{tabular}{lcccc}
\toprule
Domain & Mode & Avg IG/turn & Efficiency & Actual/Optimal turns \\
\midrule
Geography & FO & 0.80 & 80\% & $\sim$9/8 \\
Geography & PO & 0.43 & 43\% & $\sim$17/8 \\
Objects & FO & 0.73 & 73\% & $\sim$10/8 \\
Objects & PO & 0.37 & 37\% & $\sim$20/8 \\
Diseases & FO & 0.64 & 64\% & $\sim$12/8 \\
Diseases & PO & 0.35 & 35\% & $\sim$21/8 \\
\bottomrule
\end{tabular}
\end{table}

Even in the best case (geography FO), models achieve only 80\% of optimal efficiency. The disease domain is fundamentally harder due to high symptom overlap---binary search over flat symptom spaces is inefficient.

\subsection{RQ6: Prior Influence}

Full observability provides +0.36 to +0.62 win rate advantage---larger than any model or CoT manipulation.

\textbf{CoT partially compensates}: For 4B, CoT reduces the FO-PO gap by 35\% (from 0.62 to 0.40). Figure~\ref{fig:prior} shows how reasoning traces ``recover'' about 30\% of the performance lost by removing the prior.

\begin{figure}[h]
\centering
\includegraphics[width=0.7\textwidth]{fig6_prior_compensation.pdf}
\caption{CoT reduces the observability gap. Percentages show the reduction in the FO-PO win rate gap when using CoT.}
\label{fig:prior}
\end{figure}

\subsection{RQ7: Failure Analysis}

We analyzed 167,767 zero-IG turns (46\% of all turns). Figure~\ref{fig:zero_ig} shows the root cause distribution.

\begin{figure}[h]
\centering
\includegraphics[width=0.55\textwidth]{fig4_zero_ig.pdf}
\caption{Root causes of zero information gain. Critically, 43\% occur when only one candidate remains---a systematic failure to recognize convergence and assert the answer.}
\label{fig:zero_ig}
\end{figure}

\begin{table}[h]
\centering
\caption{Root causes of zero information gain turns.}
\label{tab:failures}
\small
\begin{tabular}{lcc}
\toprule
Category & Count & \% of Zero-IG \\
\midrule
Genuinely uninformative & 94,206 & 56.2\% \\
Only one candidate left & 72,298 & 43.1\% \\
Unanswerable feature & 1,263 & 0.8\% \\
Pruner error & 0 & 0.0\% \\
\bottomrule
\end{tabular}
\end{table}

\textbf{Critical finding}: 43\% of zero-IG turns occur when only one candidate remains. Models systematically fail to recognize convergence and keep asking questions instead of asserting the answer. This ``failure to guess'' is as impactful as asking uninformative questions.

\paragraph{Trajectory Clustering.}
Among 5,269 losing games, we identified three distinct failure patterns (Figure~\ref{fig:trajectories}):
\begin{itemize}
    \item 52\% \textbf{Stall}: IG drops to $\sim$0 after turn 3 and never recovers
    \item 41\% \textbf{Mixed}: Irregular, moderate IG throughout
    \item 7\% \textbf{Volatile}: High IG on turn 1 ($\sim$3.6 bits), then crash
\end{itemize}

\begin{figure}[h]
\centering
\includegraphics[width=0.95\textwidth]{fig3_trajectories.pdf}
\caption{Failure trajectory clusters. The dashed line shows optimal IG (1 bit/turn). Stall is the dominant failure mode (52\%), where models get trapped in zero-IG loops.}
\label{fig:trajectories}
\end{figure}

FO losses are 90\% stalls---when models can see candidates and still lose, it's almost always from getting stuck in uninformative questioning loops.

\section{Discussion}
\label{sec:discussion}

\paragraph{The 4B Threshold.}
Our results identify $\sim$4B parameters as a critical threshold. Below this, CoT is harmful---the model lacks capacity to leverage explicit reasoning effectively. Above it, CoT provides consistent benefits that plateau around 8B. This has practical implications: for information-seeking agents, investing in reasoning (via CoT prompting or thinking models) provides the best cost-benefit ratio at the 4-8B scale.

\paragraph{World Knowledge as Implicit Prior.}
The 235B model's shift to 65\% feature-based questions in PO mode reveals a qualitative change in strategy. Without explicit candidate metadata, large models fall back on world knowledge to generate discriminative queries (``Is it coastal?'', ``Is it chronic?''). This approximates Bayesian reasoning through implicit world models---the closest thing to emergent optimal questioning in our data.

\paragraph{The Failure-to-Guess Problem.}
Our most surprising finding is that 43\% of zero-IG turns occur with only one candidate remaining. Models are trained to ask questions, not to recognize when to stop. This suggests a missing training signal: current RLHF and preference optimization focus on question quality, not convergence detection. Future work should explicitly train models to assert answers when the hypothesis space has collapsed.

\paragraph{Domain Structure Matters.}
The dramatic difficulty gap between domains (geography: 11.8, diseases: 19.1 average difficulty, Figure~\ref{fig:difficulty}) reflects the discriminative power of available attributes. Geography has rich, hierarchical metadata (region/country/state) enabling efficient binary search. Diseases have flat, overlapping symptom lists where individual symptoms rarely eliminate more than a few candidates. This motivates domain-specific questioning strategies beyond simple binary search.

\begin{figure}[h]
\centering
\includegraphics[width=0.65\textwidth]{fig8_domain_difficulty.pdf}
\caption{Target difficulty by domain. The dashed line shows optimal (8 turns). Diseases are dramatically harder due to overlapping symptom spaces.}
\label{fig:difficulty}
\end{figure}

\section{Conclusion}

We introduced InfoGainME, a framework for evaluating LLMs' information-seeking through multi-turn dialogue games with precise information gain measurement. Our analysis of 25,513 games reveals that:

\begin{enumerate}
    \item Reasoning traces provide substantial benefits for models $\geq$4B, with the largest effect in partially observable settings.
    \item There exists a minimum scale ($\sim$4B) for effective reasoning-augmented information seeking.
    \item Large models shift from attribute-based to feature-based questioning, leveraging world knowledge as an implicit prior.
    \item Nearly half of all questions produce zero information gain, with a systematic failure to recognize convergence and assert answers.
    \item CoT eliminates redundancy, functioning as implicit dialogue memory.
\end{enumerate}

These findings provide the first rigorous characterization of how reasoning models navigate uncertainty and suggest concrete directions for improving LLM information-seeking: explicit convergence detection training, domain-aware questioning strategies, and scale-appropriate reasoning augmentation.

\bibliography{references}
\bibliographystyle{plainnat}

\end{document}
